\subsection{Vectorial faces and Tits cone}
As in the reductive case, define the fundamental chamber as $C_f^v:=\{v\in\jepPubBbb{A}\mid\forall\ i\in I,\ \alpha_i(v)>0\}$. For any subset $J$ in $I$, set
\[
F^v(J)=\{v\in\jepPubBbb{A}\mid\forall\ j\in J,\ \alpha_j(v)=0\text{ and }\forall\ i\in I\smallsetminus J,\ \alpha_i(v)>0\};
\]
then the closure $\overline{C_f^v}$ of $C_f^v$ is exactly the union of all $F^v(J)$'s for $J\subset I$. The positive vectorial faces (resp. negative vectorial faces) are defined as the sets $w\cdot F^v(J)$ (resp. $-w\cdot F^v(J)$) for $w\in W^v$ and $J\subset I$, and a vectorial face is either a positive vectorial face or a negative one. A positive chamber (resp. negative chamber) is a cone of the form $w\cdot C_f^v$ (resp. $-w\cdot C_f^v$) for some $w\in W^v$. Note that for any $x\in C_f^v$ and any $w\in W^v$, we have $(w\cdot x=1)\Rightarrow(w=1)$, which ensures that the action of $w\in W^v$ on the set of positive chambers is simply transitive.

Let $\jepPubScript{T}:=\bigcup_{w\in W^v}w\cdot\overline{C_f^v}$ be the Tits cone and $-\jepPubScript{T}$ be the negative cone. We can use it to define a $W^v$-invariant pre-order $\leqslant$ on $\jepPubBbb{A}$ as follows:
\[
\forall\ (x,y)\in\jepPubBbb{A}^2,\qquad x\leqslant y\Longleftrightarrow y-x\in\jepPubScript{T}.
\]
We also set $Y^+:=\jepPubScript{T}\cap Y$ and $Y^{++}:=Y\cap\overline{C_f^v}$. We can now recall the following simple but very useful result~\cite[Lem.~2.4 a)]{jep88refGR14}.
\begin{lemma}\label{jep88localLemma2_1}
For any $\lambda\in Y^{++}$ and any $w\in W^v$, we have $w\cdot\lambda\leqslant_{Q^{\scriptscriptstyle\vee}}\lambda$.
\end{lemma}
